\documentclass[12pt]{article}
\usepackage[utf8]{inputenc}
\usepackage{amsmath}
\usepackage{graphicx}
\usepackage{natbib}
\usepackage{hyperref}
\usepackage{lipsum} % For dummy text, you can remove this in your final document

\title{CRISPR-Cas9 in Gene Therapy for Genetic Disorders: Advancements, Clinical Trials, Ethical Considerations, and Future Directions}
\author{Your Name}
\date{\today}

\begin{document}

\maketitle

\begin{abstract}
CRISPR-Cas9 technology has revolutionized gene therapy by offering precise tools for editing the human genome. This paper reviews recent advancements in CRISPR-Cas9 applications for treating genetic disorders, discusses ongoing clinical trials, examines ethical considerations, and outlines future directions in the field of gene therapy.
\end{abstract}

\section{Introduction}
Gene therapy holds promise for treating genetic disorders by correcting disease-causing mutations at the DNA level. CRISPR-Cas9, a powerful genome editing tool, has emerged as a transformative technology in the field of molecular medicine. This paper explores the role of CRISPR-Cas9 in gene therapy, focusing on its advancements, clinical applications, ethical implications, and future prospects.

\section{Gene Therapy Overview}
\subsection{CRISPR-Cas9 Technology}
CRISPR-Cas9 enables targeted modifications of genetic sequences with unprecedented precision, facilitating therapeutic interventions for various genetic diseases \citep{jinek2012programmable, doudna2014new}.

\subsection{Applications in Genetic Disorders}
CRISPR-Cas9 has been applied to correct mutations in diseases such as cystic fibrosis, sickle cell anemia, and Duchenne muscular dystrophy, among others \citep{long2016correction, xu2019precise}.

\section{Clinical Trials}
\subsection{Current Trials and Applications}
Overview of ongoing clinical trials using CRISPR-Cas9 for gene therapy, assessing safety, efficacy, and long-term outcomes in patients \citep{maeder2013robust, liang2015targeted}.

\subsection{Case Studies}
Case studies highlighting successful applications of CRISPR-Cas9 in clinical settings, including challenges and breakthroughs in treatment outcomes \citep{tebas2014gene, tabebordbar2016efficient}.

\section{Ethical Considerations}
\subsection{Ethical Issues in Genome Editing}
Discussion on ethical concerns surrounding CRISPR-Cas9, including potential off-target effects, germline editing, and equitable access to gene therapies \citep{savulescu2015new, lanphier2015don}.

\subsection{Regulatory Frameworks}
Overview of regulatory frameworks governing gene therapy and genome editing technologies, ensuring ethical guidelines and patient safety \citep{regalado2015who, hsu2014regulatory}.

\section{Future Directions}
\subsection{Technological Advancements}
Potential advancements in CRISPR-Cas9 technology, such as base editing, prime editing, and enhanced delivery methods for broader therapeutic applications \citep{ansell2017genome, anzalone2019search}.

\subsection{Clinical and Translational Research}
Future directions in clinical research, including scaling up gene therapies, addressing technical challenges, and expanding treatment accessibility globally \citep{cong2013multiplex, nelson2015genome}.

\section*{References}
\bibliographystyle{plainnat}
\bibliography{prompt14_35}

\end{document}
